\subsection{Aplicación a los cambios de base en los preesquemas algebraicos.}

\begin{proposition}[6.7.1]
\label{IV.2.s6.7.1-fr}
Sean $k$ un cuerpo, $X$ un $k$-preesquema localmente noetheriano y
$\mathscr F$ un $\mathcal O_X$-módulo coherente. Sea $k'$ una extensión
de $k$; pongamos $X'=X\otimes_k k'$, $\mathscr F'=\mathscr F\otimes_k k'$
y sea $p:X'\to X$ la proyección canónica. Se supone, bien que
$X$ es localmente de tipo finito sobre $k$, bien que $k'$ es una extensión
de tipo finito de $k$, de modo que en ambos casos $X'$ es localmente
noetheriano. Sean $x'$ un punto de $X'$ y $x=p(x')$. Entonces:
\begin{enumerate}
\item[\textup{(i)}] Se tiene $\operatorname{coprof}(\mathscr F_x)=
\operatorname{coprof}(\mathscr F'_{x'})$; en particular, para que
$\mathscr F_x$ sea un $\mathcal O_x$-módulo de Cohen-Macaulay, es
necesario y suficiente que $\mathscr F'_{x'}$ sea un
$\mathcal O_{x'}$-módulo de Cohen-Macaulay.
\item[\textup{(ii)}] Para que $\mathscr F$ verifique la propiedad $(S_n)$ en el punto
$x$, es necesario y suficiente que $\mathscr F'$ verifique $(S_n)$ en el
punto $x'$.
\end{enumerate}
\end{proposition}

\oldpage[IV-2]{146}

\begin{proof}
Se sabe que $p$ es fielmente plano (2.2.13); las dos afirmaciones son, por
tanto, consecuencias respectivas de (6.3.2) y (6.4.2), con tal de que se
demuestre que $p^{-1}(x)=\operatorname{Spec}(\boldsymbol{k}(x)\otimes_k k')$ es un
\emph{preesquema de Cohen-Macaulay}; como uno de los dos cuerpos $\boldsymbol{k}(x)$, $k'$ es
una extensión de tipo finito de $k$ por hipótesis la cuestión se reduce a demostrar
el
\end{proof}

\begin{lemma}[6.7.1.1]
\label{IV.2.s6.7.1.1-fr}
Sean $K$, $L$ dos extensiones de un cuerpo $k$, una de las cuales es de tipo
finito (de modo que el anillo $K\otimes_k L$ es noetheriano). Entonces
$K\otimes_k L$ es un anillo de Cohen-Macaulay.
\end{lemma}

\begin{proof}
Supongamos por ejemplo que $L$ es una extensión de tipo finito de $k$,
de modo que $L$ es una extensión finita de una extensión puramente
trascendente $k(\mathbf t)$ de $k$ (siendo $\mathbf t$ un sistema
$(t_i)_{1\leqslant i\leqslant m}$ de indeterminadas). Si se pone
$A=K\otimes_k k(\mathbf t)$, $B=K\otimes_k L$, $B$ es un $A$-módulo
plano $(\mathbf{0}_{\mathrm I},6.2.1)$ y de tipo finito; el morfismo
$h:\operatorname{Spec}(B)\to\operatorname{Spec}(A)$ es, por tanto, finito,
y como todo preesquema artiniano es de Cohen-Macaulay, para demostrar que
$B$ es un anillo de Cohen-Macaulay basta demostrar que $A$ es un anillo de
Cohen-Macaulay (6.3.5). Ahora bien, si $S$ es el conjunto de los elementos
$\neq 0$ de $k[\mathbf t]$, se tiene $A=S^{-1}A'$, donde
$A'=K\otimes_k k[\mathbf t]=K[\mathbf t]$; en virtud de
(\textbf 0, 16.5.13), basta demostrar que $A'$ es un anillo de Cohen-Macaulay; pero
esto resulta de que $A'$ es \emph{regular} (\textbf 0, 17.3.7).
\end{proof}

\begin{corollary}[6.7.2]
\EGAReaderTarget{ega4.c04.cor.6-7-2}{IV.2.s6.7.2-fr}
Para que $\mathcal O_x$ sea un anillo de Cohen-Macaulay \textup{(resp.} para
que $X$ verifique la propiedad $(S_n)$ en el punto $x$), es necesario y
suficiente que $\mathcal O_{x'}$ sea un anillo de Cohen-Macaulay \textup{(resp.} que
$X'$ verifique $(S_n)$ en el punto $x'$).
\end{corollary}

\begin{corollary}[6.7.3]
\label{IV.2.s6.7.3-fr}
Sean $X$, $Y$ dos $k$-preesquemas localmente noetherianos, siendo al menos
uno de ellos localmente de tipo finito sobre $k$. Sea $\mathscr F$ \textup{(resp.}
$\mathscr G$) un $\mathcal O_X$-módulo coherente \textup{(resp.} un
$\mathcal O_Y$-módulo coherente). Si $\mathscr F$ y $\mathscr G$
verifican la propiedad $(S_n)$, lo mismo ocurre con
$\mathscr F\otimes_k\mathscr G$.
\end{corollary}

\begin{proof}
La hipótesis implica que $X\times_kY$ es localmente noetheriano; denotemos
por $p:X\times_kY\to X$ y $q:X\times_kY\to Y$ las proyecciones canónicas,
que son morfismos planos. Supongamos por ejemplo que $X$ es localmente de
tipo finito sobre $k$; se puede escribir
$\mathscr F\otimes_k\mathscr G=p^*(\mathscr F)\otimes_k\mathscr G$
y puesto que $\mathscr F$ es plano relativamente al morfismo estructural
$X\to\operatorname{Spec}(k)$, $p^*(\mathscr F)$ es $q$-plano (2.1.4).
Apliquemos el criterio (6.4.1, (ii)) al morfismo $q$; basta ver que, para
todo $y\in Y$, $p^*(\mathscr F)\otimes_{\mathcal O_Y}\boldsymbol{k}(y)$ posee la
propiedad $(S_n)$; pero
$p^*(\mathscr F)\otimes_{\mathcal O_Y}\boldsymbol{k}(y)=\mathscr F\otimes_k \boldsymbol{k}(y)$,
y como $X$ es localmente de tipo finito sobre $k$, la conclusión resulta
de (6.7.1, (ii)).
\end{proof}
